\documentclass[conference]{IEEEtran}
\usepackage{cite}
\usepackage{amsmath,amssymb}
\usepackage{graphicx}
\usepackage{booktabs}
\usepackage{url}
\usepackage[hidelinks]{hyperref}

\graphicspath{{figures/}}

\begin{document}

\title{Causal Consistency as a Defense: Detecting Provenance Graph Poisoning}

\author{\IEEEauthorblockN{N.~Manoj}
\IEEEauthorblockA{Independent Researcher\\
Email: nehamanoj1105@users.noreply.github.com}}

\maketitle

\begin{abstract}
Provenance-based intrusion detection operates on the assumption that the
provenance graph it consumes faithfully records what happened on the host. If an
adversary can tamper with the stored graph after collection, a downstream
detector reasons over a fiction. We study whether the causal structure of a
provenance graph is self-validating: whether a fixed set of deterministic
consistency rules can expose post-collection tampering without any training.
We evaluate a 15-rule engine against a GraphSAGE baseline under four tampering
operations on synthetic graphs and on real DARPA Transparent Computing
Engagement~3 Theia provenance, and we examine both under adversarial mimicry.
On controlled synthetic poisoning the rule engine reaches F1
$0.805\pm0.063$ over ten seeds, against F1 $0.322\pm0.194$ for a leakage-free
GraphSAGE, a paired difference of $+0.482$ (Wilcoxon $p=0.002$). With a genuine
continuous anomaly score the rule engine attains ROC-AUC $0.953$, and its
accuracy is unchanged as mimicry strength grows from none to heavy. On real
Theia graphs with synthetic post-collection poisoning the rule engine reaches
F1 $0.83$--$0.85$ where GraphSAGE is near chance. We also report a methodological
finding: the naive training procedure commonly used for such baselines leaks the
evaluation labels into threshold selection and inflates GraphSAGE F1 from
$0.246$ to $0.588$. All results, protocols, and the corrected evaluation are
released for independent verification.
\end{abstract}

\begin{IEEEkeywords}
provenance, intrusion detection, graph poisoning, causal consistency,
GraphSAGE, evaluation methodology
\end{IEEEkeywords}

\section{Introduction}
End-to-end provenance captures the causal history of a system as a directed
graph whose nodes are processes, files, users and network endpoints and whose
edges are read, write, execute, connect, spawn and delete events. Systems such
as Unicorn, NoDoze, HOLMES, Poirot and MAGIC use this graph to detect advanced
persistent threats, and they do so under an assumption that is rarely stated:
the graph they analyze is authentic. The graph is produced by a logging pipeline
and then stored, shipped, and indexed. Every stage after kernel-level capture is
a potential point of compromise. An adversary who can edit the stored graph can
delete the events that incriminate them, insert events that exonerate them, or
reorder events to break the causal chain a detector relies on.

This paper asks a narrow, testable question: to what extent is causal
consistency in a provenance graph sufficient to reveal such tampering? We do not
propose a new learning architecture. We take the invariants that a provenance
graph must satisfy by construction---an execute edge must point at a file, a
read edge must point at a file, a child's spawn must follow its parent's spawn,
event sequence numbers must be monotone within a process---and assemble them
into a deterministic rule engine. We then compare that engine to a graph neural
network baseline, GraphSAGE, on the same poisoned instances.

The comparison is deliberately adversarial to our own method. A learned detector
should, in principle, capture correlations that hand-written rules cannot. What
we find is that the comparison is dominated not by modeling capacity but by
evaluation design. A GraphSAGE classifier trained and thresholded in the
conventional way reaches an F1 that is largely an artifact of selecting its
decision threshold on the very labels it is scored against. Once the split is
made honest, the learned baseline loses most of its apparent advantage, while
the training-free rule engine is unaffected because it has no split to leak.

Our contributions are:
\begin{itemize}
\item a deterministic, training-free detector for post-collection provenance
tampering built from fifteen causal-consistency rules (Section~\ref{sec:method});
\item a rigorous evaluation protocol with disjoint train, validation and test
edges and validation-only threshold selection, and the demonstration that the
conventional protocol inflates GraphSAGE F1 by more than $0.34$ absolute
(Section~\ref{sec:eval});
\item a ten-seed comparison on synthetic poisoning and an evaluation on real
Theia provenance with synthetic post-collection poisoning, with paired
significance tests (Sections~\ref{sec:results}--\ref{sec:robust});
\item an audit showing that a reported ROC-AUC of $1.000$ for the rule engine
is not a computed quantity, and the substitution of a genuine continuous score
that yields ROC-AUC $0.953$ (Section~\ref{sec:robust}).
\end{itemize}

\section{Related Work}
Provenance-based detection spans specification-based, statistical, and
learning-based systems. Specification systems such as HOLMES and Poirot encode
expert knowledge of attack tactics and are precise but costly to maintain.
Statistical systems such as NoDoze and ProvDetector score the rarity of events
or paths and are prone to flagging unusual-but-benign activity. Learning-based
systems---Unicorn, FLASH, MAGIC and ShadeWatcher---learn representations of
benign behavior from the graph and detect deviations; GraphSAGE is a common
backbone for these methods.

The vulnerability we study is adjacent to adversarial mimicry against graph
detectors, in which an attacker interleaves benign events with malicious ones to
dilute the anomaly signal. It is also adjacent to the broader integrity problem
for provenance: if the graph is editable, every downstream conclusion is
conditional on that graph. Our work differs from prior detection papers in
framing: rather than detecting the attacker's behavior inside a trusted graph,
we verify the graph itself, and we are explicit that the rules are checking
structural and causal well-formedness rather than learning a behavioral model.

\section{Threat Model}
We assume the attacker can modify the stored provenance graph after collection
but cannot alter the detector, disable the parser, or compromise the kernel that
produced the events. Four operations are in scope: edge deletion (removing a
recorded event), edge insertion (fabricating an event that did not occur),
edge reordering (changing timestamps so that events no longer respect causal
order), and dependency forgery (re-attributing an event to a different source or
target). The attacker is gray-box: they know the graph schema and the general
class of detector, but not its internal parameters.

The detection target is the set of modified edges in the final graph. We are
careful about what this means for deletion. A deleted edge no longer exists in
the graph, so an edge-level detector cannot mark it: deletion is only detectable
insofar as removing an event breaks a higher-level invariant elsewhere (for
example, a process left without its spawn record, or a gap in an event
sequence). We therefore report deletion separately and do not fold it into the
edge-level metrics as if it were a scored positive.

\section{Provenance Representation}
A provenance graph $G=(V,E)$ has typed nodes and edges. Node types are process,
file, user and network endpoint. Edge types are read, write, execute, connect,
spawn and delete. Each edge carries a source identifier, a target identifier, a
type, and a timestamp; each event carries a sequence number within its process.
The schema is the common data model used by the DARPA Transparent Computing
program, and our parser maps its Avro records into this representation.

\section{Proposed Detection Framework}
\label{sec:method}
The detector is a rule engine that evaluates fifteen invariants over $G$. Each
rule returns a set of violations, each attributed to an edge or a node. The
rules fall into three groups.

\textbf{Structural rules} check well-formedness: no duplicate edge tuples, no
duplicate event identifiers, spawn edges must target process nodes, execute and
read/write/delete edges must target file nodes, connect edges must target
network endpoints, no self-referential edges, no edges referencing missing
nodes, and no active process lacking a parent spawn edge.

\textbf{Temporal rules} check ordering: non-negative timestamps, child spawn
after parent spawn, process activity after process spawn, no gaps in event
sequence numbers, and no timestamp inversions within a process stream.

\textbf{Semantic rules} check lineage: processes that appear to act but were
never spawned by any observed parent.

The engine is deterministic and requires no training data. For ranking and for
ROC-AUC we assign each edge a continuous anomaly score: every violation
contributes a severity weight (high $3$, medium $2$, low $1$) to the edge it is
attributed to, either directly by edge identifier or, for node-attributed
violations, to the incident edges. The binary decision is score greater than
zero. This score is a genuine ranking signal, and we use it rather than any
post-hoc construction when reporting AUC.

\section{GraphSAGE Baseline}
We include GraphSAGE as a representative learned detector. Nodes are featurized
by a one-hot node type and degree; edges by a one-hot edge type and a normalized
timestamp. A two-layer GraphSAGE encoder produces edge embeddings that a
classifier maps to a tampering probability. The baseline exists to test whether
learning can outperform fixed invariants; it is not tuned to win, and our main
methodological result concerns how it is evaluated.

\section{Experimental Methodology}
\label{sec:eval}
\subsection{Datasets}
We use two sources. The first is a synthetic provenance generator that produces
a graph of thirty processes, forty files and ten network entities; this gives
full control over ground truth. The second is real provenance from DARPA
Transparent Computing Engagement~3 (Theia). Two of the four benchmark scenarios
were obtained and parsed: scenario~3 (20{,}157 nodes, 297{,}777 edges) and
scenario~5m (34{,}835 nodes, 464{,}858 edges). Scenarios~1r and~6r could not be
downloaded because the public mirror enforced a download quota; we report
results for the two obtained scenarios and state the missing coverage rather
than estimating it.

\subsection{Poisoning}
On each graph we inject five deletions, five insertions, five reorderings and
five dependency forgeries, for twenty events per seed. The injector operates on
a deep copy of the clean graph, records the original and modified state of each
event, and asserts that the clean graph is unmodified afterwards. Twenty events
produce fifteen scored positives, because the five deletions leave no
final-graph record. Detection metrics are computed over the edges present in the
final graph; the deleted edges are reported as undetectable under edge matching
and excluded. This is the \emph{final-graph universe}. We note that an
alternative convention, which adds the identifiers of deleted edges to the
scoring universe, raises F1 from $0.805$ to $0.845$; we report the
final-graph universe and state the assumption.

\subsection{Leakage-free protocol}
For the learned baseline, edges are partitioned into disjoint training,
validation and test sets. Model weights are fit on training edges only, the
decision threshold is selected on validation edges only, and every reported
metric is computed on the untouched test set. Message passing for test-node
features uses training and validation edges only, so the test edges neither
influence the weights nor the threshold. For the rule engine there is nothing to
leak, since it has no fitted parameters and no threshold search.

The conventional protocol found in much of the reference code does none of
this: it fits on all edges and then chooses the threshold that maximizes F1 on
the same labels it reports. We measure the effect of this directly.

\subsection{Seeds and statistics}
All synthetic, mimicry, ablation and real-Theia experiments run over ten seeds
($1,7,13,21,42,99,123,256,512,1024$); generalization runs over five seeds
($1,7,13,21,42$). We report the mean, sample standard deviation, median, minimum,
maximum and a 95\% confidence interval (Student-$t$). Comparisons between the
rule engine and GraphSAGE use the paired Wilcoxon signed-rank test on identical
test instances, with Cohen's $d_z$.

\section{Main Results}
\label{sec:results}
\subsection{Synthetic controlled poisoning}
Table~\ref{tab:main} reports the ten-seed comparison. The rule engine attains F1
$0.805\pm0.063$ with precision $0.715$ and recall $0.933$, and a high specificity
of $0.964$; its MCC is $0.795$. GraphSAGE, under the leakage-free protocol,
attains F1 $0.322\pm0.194$ with precision $0.342$ and recall $0.502$. The learned
baseline is not competitive once its evaluation is honest, and its variance
across seeds is far larger.

The rule engine's ROC-AUC is $0.953\pm0.031$ and its PR-AUC is $0.740$; these are
computed from the continuous severity-weighted score described in
Section~\ref{sec:method}. GraphSAGE reaches ROC-AUC $0.829$. Figure~\ref{fig:main}
shows the comparison, and Figure~\ref{fig:rocpr} shows the rule engine's ROC and
precision--recall curves for a representative seed.

\subsection{Why the comparison is dominated by evaluation design}
The most consequential number in this paper is not a detection score. When the
GraphSAGE baseline is trained and thresholded in the conventional way---weights
on all edges, threshold maximizing F1 on the scored labels---it reaches F1
$0.588$ on the same seed where the leakage-free protocol yields $0.246$
(Figure~\ref{fig:leakage}). The gap is not model quality; it is the threshold
being chosen with knowledge of the answers. Reports that compare a learned
baseline to a training-free detector without a held-out threshold search are
therefore comparing the baseline's leaked threshold to the detector's fixed
rule.

\subsection{Real Theia provenance with synthetic poisoning}
Table~\ref{tab:realpoison} reports the real-data experiments. On scenario~3 the
rule engine reaches F1 $0.835\pm0.046$ and on scenario~5m F1 $0.850\pm0.060$;
GraphSAGE is near chance on both (F1 $0.034$ and $0.010$). Two caveats are
essential. First, the attacks here are synthetic post-collection injections into
real provenance; they are not real DARPA attack labels, and we do not claim
detection of naturally occurring attacks. Second, we truncate the real graphs to
their first 50{,}000 edges, as the original loader does; the truncation is
stated, not hidden. On the unmodified real graphs the rule engine flags only
$19$ of $50{,}000$ edges in scenario~3 and $7$ of $50{,}000$ in scenario~5m, a
false-positive rate below $0.04\%$.

\section{Robustness and Ablation}
\label{sec:robust}
\subsection{Mimicry}
We generate camouflage edges that preserve the degree, type and timestamp
distributions of the surrounding graph and, critically, do not themselves
violate any invariant. Table~\ref{tab:mimicry} and Figure~\ref{fig:mimicry}
show the result: the rule engine's F1 is $0.805$ at every strength from none to
heavy, unchanged, because camouflage that satisfies the invariants is
computationally indistinguishable from benign behavior under a consistency
check. GraphSAGE moves between F1 $0.304$ and $0.542$ and stays well below the
rule engine throughout. We emphasize a specification detail: mimicry noise in
this experiment is constructed to be genuinely consistent. If camouflage is
generated so that it violates the rules it is meant to evade---for instance by
giving noise edges inconsistent types or timestamps---the rule engine flags the
noise, and any apparent precision collapse is an artifact of the generator, not
a property of the attack.

\subsection{Rule ablation}
Table~\ref{tab:ablation} reports the ablation over ten seeds. The structural
group carries the detector: removing it drops F1 from $0.805$ to $0.328$.
Removing the temporal group costs little ($0.805$ to $0.798$), and the
leave-one-rule-out analysis identifies the read/write, network and spawn
consistency rules as the individually load-bearing ones (removing read/write
alone drops F1 to $0.517$). Two rules, the unspawned-process and timestamp
rules, are inert on synthetic data: removing either changes nothing. We report
this rather than presenting all fifteen rules as equally necessary; their value
is expected to appear on real logs with non-sequential identifiers, which is
future work.

\subsection{Cross-dataset generalization}
The rule engine requires no retraining across graphs. Applied unchanged to a
real Theia graph after being defined on synthetic data, it reaches F1 $0.857$;
moving between the two real scenarios it reaches $0.837$ (Table~\ref{tab:generalization}).
We make no generalization claim for GraphSAGE: a genuine transfer across a
shared feature space is not implemented in this study, so the GraphSAGE column
in Table~\ref{tab:generalization} is a target-trained model and is labelled as such.

\section{Scalability}
Table~\ref{tab:scalability} reports runtime, memory and throughput as graph size
grows from ten thousand to one million edges, five repetitions per size.
Construction and rule-engine inference are both linear in the number of edges.
Inference takes $0.37$~s at ten thousand edges and $34.5$~s at one million, and
peak Python allocation grows from roughly $37$~MB to $971$~MB, about one
kilobyte per edge. Throughput is stable near $2.9\times10^{4}$ edges per second
across three orders of magnitude (Figure~\ref{fig:scalability}), which is what a
single linear pass over the events should give. We report GraphSAGE training and
inference separately from rule-engine inference and do not compare a training
cost against an inference cost.

\section{Limitations}
Three limitations bound our claims. First, the real-data evaluation covers two
of four Theia scenarios; scenarios~1r and~6r were unavailable, and no results
are estimated for them. Second, all detection of DARPA provenance in this paper
involves synthetic post-collection poisoning, not real attacks; the real-data
experiments measure whether the detector works on real graph structure, not
whether it catches real intrusions. Third, deletion is only partially
detectable: a deleted edge cannot be scored by edge matching, so our metrics
exclude it, and we do not claim a deletion detection rate. Two further
limitations concern the baseline rather than the method: the GraphSAGE feature
set is minimal, and no shared-feature cross-graph transfer is implemented.

\section{Conclusion}
A provenance graph carries enough internal causal structure to reveal that it
has been edited, and a small deterministic rule set exploits that structure
without training. Across ten seeds on controlled poisoning the rule engine
outperforms a leakage-free GraphSAGE by $0.48$ F1 with a large paired effect,
its accuracy is invariant to invariant-preserving mimicry, and it scales
linearly to a million-edge graphs. The comparison also carries a methodological
lesson: the size of a baseline's reported advantage can be an artifact of
selecting its threshold on the evaluation labels, and removing that leak removes
most of the advantage. We release the complete protocol, raw per-seed results
and the corrected evaluation so that every number here can be re-derived.

\input{tables/table3_main_detection}
\input{tables/table9_real_data_poisoning}
\input{tables/table6_mimicry}
\input{tables/table7_rule_ablation}
\input{tables/table5_cross_dataset}
\input{tables/table8_scalability}

\begin{figure}[t]
\centering
\includegraphics[width=0.95\columnwidth]{fig_main_performance.pdf}
\caption{Synthetic controlled poisoning, ten seeds: Rule Engine versus
leakage-free GraphSAGE (mean and 95\% CI).}
\label{fig:main}
\end{figure}

\begin{figure}[t]
\centering
\includegraphics[width=0.95\columnwidth]{fig_roc_pr.pdf}
\caption{Rule Engine ROC and precision--recall curves from a genuine continuous
anomaly score (synthetic seed 42).}
\label{fig:rocpr}
\end{figure}

\begin{figure}[t]
\centering
\includegraphics[width=0.8\columnwidth]{fig_leakage.pdf}
\caption{Effect of evaluation leakage: the same GraphSAGE seed reaches F1
$0.588$ when the threshold is tuned on the scored labels and $0.246$ when the
threshold is selected on held-out validation edges.}
\label{fig:leakage}
\end{figure}

\begin{figure}[t]
\centering
\includegraphics[width=0.95\columnwidth]{fig_mimicry.pdf}
\caption{Robustness to invariant-preserving mimicry (ten seeds). The Rule
Engine's F1 is unchanged; mimicry noise introduces zero rule self-violations.}
\label{fig:mimicry}
\end{figure}

\begin{figure}[t]
\centering
\includegraphics[width=0.95\columnwidth]{fig_scalability.pdf}
\caption{Scalability, five repetitions per size. Construction and inference are
linear in edge count and throughput is stable.}
\label{fig:scalability}
\end{figure}

\begin{thebibliography}{99}
\bibitem{han2020unicorn} X.~Han, T.~Pasquier, A.~Bates, J.~Mickens, and
M.~Seltzer, ``Unicorn: Runtime provenance-based detector for advanced
persistent threats,'' in \emph{NDSS}, 2020.
\bibitem{hsieh2018kobold} K.~Hsieh \emph{et al.}, ``Kobold: Evaluating
decentralized access control,'' in \emph{IEEE S\&P Workshops}, 2018.
\bibitem{jiang2023mimicry} P.~Jiang \emph{et al.}, ``Mimicry attacks against
provenance graph based intrusion detection systems,'' in \emph{NDSS}, 2023.
\bibitem{wang2020you} Q.~Wang \emph{et al.}, ``You are what you do: Hunting
stealthy malware via data provenance analysis,'' in \emph{NDSS}, 2020.
\bibitem{li2021magnet} Z.~Jia, Y.~Xiong, Y.~Nan, Y.~Zhang, J.~Zhao, and
M.~Wen, ``MAGIC: Detecting advanced persistent threats via masked graph
representation learning,'' in \emph{USENIX Security}, 2024.
\bibitem{chen2021flash} M.~A.~Rehman and S.~Ahmadi, ``FLASH: A comprehensive
approach to intrusion detection via provenance graph representation learning,''
in \emph{IEEE S\&P}, 2022.
\bibitem{king2003backtracking} S.~T.~King and P.~M.~Chen, ``Backtracking
intrusions,'' in \emph{ACM SOSP}, 2003.
\bibitem{hossain2017sleuth} M.~N.~Hossain \emph{et al.}, ``SLEUTH: Real-time
attack scenario reconstruction from COTS audit data,'' in \emph{USENIX
Security}, 2017.
\bibitem{milajerdi2019holmes} S.~M.~Milajerdi, R.~Gjomemo, B.~Eshete,
R.~Sekar, and V.~N.~Venkatakrishnan, ``HOLMES: Real-time APT detection through
correlation of suspicious information flows,'' in \emph{IEEE S\&P}, 2019.
\bibitem{peirce2021poirot} S.~M.~Milajerdi, B.~Eshete, R.~Gjomemo, and
V.~N.~Venkatakrishnan, ``Poirot: Aligning attack behavior with kernel audit
records for cyber threat hunting,'' in \emph{ACM CCS}, 2019.
\bibitem{hasan2021nodoze} W.~U.~Hassan, M.~A.~Noureddine, P.~Datta, and
A.~Bates, ``NoDoze: Combatting threat alert fatigue with automated provenance
triage,'' in \emph{NDSS}, 2019.
\bibitem{hamilton2017graphsage} W.~L.~Hamilton, R.~Ying, and J.~Leskovec,
``Inductive representation learning on large graphs,'' in \emph{NeurIPS}, 2017.
\bibitem{darpa2020tc} DARPA I2O, ``Transparent Computing Engagement 3 data
release,'' 2020. [Online]. Available:
\url{https://github.com/darpa-i2o/Transparent-Computing}
\bibitem{pasquier2018practical} T.~Pasquier \emph{et al.}, ``Practical
whole-system provenance capture,'' in \emph{ACM SoCC}, 2017.
\end{thebibliography}

\end{document}
